With the same terminal weight and 160 updates, dense plus phase feedback reached 65.3\%, compared with 61.7\% for terminal feedback alone. On a separate task starting at 2\% success, dense feedback reached 18\% after 20 updates and terminal feedback reached 11\%. Both runs used the same group filter. These observations show higher measured complete-task success with dense feedback in these comparisons; they do not isolate an improvement in individual skills. The 110-update comparison cannot isolate the effect of continuous progress because the phase weight also changes.

Applying only a binary terminal reward to the CDPR also produced successful placements. The terminal-only and phased-reward runs used different starting states, supervision, and success criteria, so the current experiments do not establish a causal gain from switching to a phased reward. Such a comparison requires the same initializer, resets, termination rule, and interaction budget.

All robot results are from simulation. The CDPR repeats measure rollout variation for one trained policy lineage on development scenes used for model selection. The LIBERO comparisons report one endpoint per run. Independent training runs, the reserved final-test scenes, and physical execution remain untested. The contribution of scripted composition data also remains to be isolated. The results support the use of skill acquisition, recovery, and explicit transition evaluation, while the weak WidowX pickup result shows that separately learned placement cannot compensate for an unreliable earlier skill.
